The definitions of "URI-reference", "absolute-URI", "relative-part", "authority", "port", "host", "path-abempty", "segment", and "query" 
are adopted from the URI generic syntax.

An "absolute-path" rule is defined for protocol elements that can contain a non-empty path component.
(This rule differs slightly from the path-abempty rule of RFC 3986, which allows for an empty path, 
and path-absolute rule, which does not allow paths that begin with "//".) A "partial-URI" rule is defined for 
protocol elements that can contain a relative URI but not a fragment component.

  URI-reference = <URI-reference, see [URI], Section 4.1>
  absolute-URI  = <absolute-URI, see [URI], Section 4.3>
  relative-part = <relative-part, see [URI], Section 4.2>
  authority     = <authority, see [URI], Section 3.2>
  uri-host      = <host, see [URI], Section 3.2.2>
  port          = <port, see [URI], Section 3.2.3>
  path-abempty  = <path-abempty, see [URI], Section 3.3>
  segment       = <segment, see [URI], Section 3.3>
  query         = <query, see [URI], Section 3.4>

  absolute-path = 1*( "/" segment )
  partial-URI   = relative-part [ "?" query ]

  hier-part		= "//" authority path-abempty
				/ path-absolute
                / path-rootless
                / path-empty

IP ADDDRESS
===========
IPv6address   =                            6( h16 ":" ) ls32
                 /                       "::" 5( h16 ":" ) ls32
                 / [               h16 ] "::" 4( h16 ":" ) ls32
                 / [ *1( h16 ":" ) h16 ] "::" 3( h16 ":" ) ls32
                 / [ *2( h16 ":" ) h16 ] "::" 2( h16 ":" ) ls32
                 / [ *3( h16 ":" ) h16 ] "::"    h16 ":"   ls32
                 / [ *4( h16 ":" ) h16 ] "::"              ls32
                 / [ *5( h16 ":" ) h16 ] "::"              h16
                 / [ *6( h16 ":" ) h16 ] "::"

   h16           = 1*4HEXDIG
   ls32          = ( h16 ":" h16 ) / IPv4address
   IPv4address   = dec-octet "." dec-octet "." dec-octet "." dec-octet

   dec-octet     = DIGIT                 ; 0-9
                 / %x31-39 DIGIT         ; 10-99
                 / "1" 2DIGIT            ; 100-199
                 / "2" %x30-34 DIGIT     ; 200-249
                 / "25" %x30-35          ; 250-255

URI-reference = URI / relative-ref
===================================
1: URI-reference = scheme ":" hier-part [ "?" query ] [ "#" fragment ]
		 
2: URI-reference = relative-ref  = relative-part [ "?" query ] [ "#" fragment ]
		 	   
	relative-part = "//" authority path-abempty
    	            / path-absolute
        	        / path-noscheme
            		/ path-empty

Each protocol element in HTTP that allows a URI reference will indicate in its ABNF production 
      
      1) whether the element allows any form of reference (URI-reference), 
      2) only a URI in absolute form (absolute-URI),
      3) only the path and optional query components (partial-URI)
      4) some combination of the above. 
  
Unless otherwise indicated, URI references are parsed relative to the target URI (Section 7.1). 

It is RECOMMENDED that all senders and recipients support, at a minimum, URIs with lengths of 8000 octets in protocol elements. 
Note that this implies some structures and on-wire representations (for example, the request line in HTTP/1.1) will necessarily be larger in some cases.

Absolute-path : defined for protocol elements that can contain a non-empty path component
  Elements that can contain a non-empty path component :





URI RULES
=========
1. A sender MUST NOT generate an "http" URI with an empty host identifier : A recipient that processes such a URI reference MUST reject it as invalid.
    i.e. URI = "http://" or "http:///abdc" 
2. If the port is equal to the default port for a scheme, the normal form is to omit the port subcomponent.When not being used as the target of an OPTIONS request, 
    an empty path component is equivalent to an absolute path of "/", so the normal form is to provide a path of "/" instead.
3. The scheme and host are case-insensitive (normally provided in lowercase); 
4. All other components are compared in a case-sensitive manner.
5. Characters other than those in the "reserved" set are equivalent to their percent-encoded octets: 
    the normal form is to not encode them (see Sections 2.1 and 2.2 of [URI]).

scheme      = ALPHA *( ALPHA / DIGIT / "+" / "-" / "." )
authority   = [ userinfo "@" ] host [ ":" port ] 
  userinfo = *( unreserved / pct-encoded / sub-delims / ":" )
  host        = IP-literal / IPv4address / reg-name
      IP-literal = "[" ( IPv6address / IPvFuture  ) "]"
      IPvFuture  = "v" 1*HEXDIG "." 1*( unreserved / sub-delims / ":" )
      reg-name    = *( unreserved / pct-encoded / sub-delims )
      port        = *DIGIT
  
path          = path-abempty    ; begins with "/" or is empty
                    / path-absolute   ; begins with "/" but not "//"
                    / path-noscheme   ; begins with a non-colon segment
                    / path-rootless   ; begins with a segment
                    / path-empty      ; zero characters
  
  path-abempty  = *( "/" segment )  [ANY NUMBER OF /PCHAR/PCHAR]
  path-absolute = "/" [ segment-nz *( "/" segment ) ]
  path-noscheme = segment-nz-nc *( "/" segment )
  path-rootless = segment-nz *( "/" segment )
  path-empty    = 0<pchar>

  segment       = *pchar
  segment-nz    = 1*pchar
  segment-nz-nc = 1*( unreserved / pct-encoded / sub-delims / "@" )
                ; non-zero-length segment without any colon ":"

  pchar         = unreserved / pct-encoded / sub-delims / ":" / "@"

query       = *( pchar / "/" / "?" )
 fragment    = *( pchar / "/" / "?" )



  URI REFERENCE
  =============
  URI-reference = URI / relative-ref
    if not a URI (i.e "recognised-scheme:") then URI-reference is a relative-reference (AKA target URI)

    relative-ref  = relative-part [ "?" query ] [ "#" fragment ]

    relative-part = "//" authority path-abempty
                  / path-absolute
                  / path-noscheme
                  / path-empty

        A relative reference that begins with two slash characters is termed a network-path reference; : such references are rarely used.  
        A relative reference that begins with a single slash character is termed an absolute-path reference.
        A relative reference that does not begin with a slash character is termed a relative-path reference.  
       
        A path segment that contains a colon character (e.g., "this:that") cannot be used as the first segment 
          of a relative-path reference, as it would be mistaken for a scheme name. 
        Such a segment must be preceded by a dot-segment (e.g., "./this:that") to make a relative-path reference

URI form
=========
request-target = origin-form
                 / absolute-form
                 / authority-form
                 / asterisk-form

asterisk-form  = "*"													e.g. OPTIONS * HTTP/1.1
absolute-form = absolute-URI										
absolute-URI 	[absolute-URI  = scheme ":" hier-part [ "?" query ]] 	e.g GET http://www.example.org/pub/WWW/TheProject.html HTTP/1.1
absolute-path 	[absolute-path = 1*( "/" segment )]						
authority-form	[uri-host ":" port]										e.g. www.example.com:80
origin-form   	[absolute-path [ "?" query ]]							e.g. /search?q=openai+http HTTP/1.1							

Target URI 
===========
The target URI is the request-target when the request-target is in absolute-form. 

target URI is [absolute-URI  = scheme ":" hier-part [ "?" query ]]